%% ---------------------------------------------------------------------------
%%  Context-Dependent Selection of Urban Routing Policies
%%  IC_TORS'26 submission manuscript.  Standalone: compile main.tex directly.
%%    pdflatex main -> bibtex main -> pdflatex main -> pdflatex main
%%  Figures are external PDFs in figures/.  Tables are inline.
%% ---------------------------------------------------------------------------
\documentclass[runningheads]{llncs}

\usepackage[T1]{fontenc}
\usepackage[utf8]{inputenc}
\usepackage{lmodern}
\usepackage{amsmath,amssymb}
\usepackage{graphicx}
\usepackage{booktabs}
\usepackage{tabularx}
\usepackage{array}
\usepackage{microtype}
\usepackage[hidelinks]{hyperref}
\usepackage{url}

\renewcommand{\UrlFont}{\ttfamily\small}
\newcommand{\CO}{CO\textsubscript{2}}
\newcolumntype{R}{>{\raggedright\arraybackslash}X}
\setlength{\tabcolsep}{4pt}

\begin{document}

\title{Context-Dependent Selection of Urban Routing Policies:\\
Regime Effects, Decision Value and Generalisation Limits}
\titlerunning{Context-Dependent Selection of Urban Routing Policies}

\author{Anonymous author(s)}
\authorrunning{Anonymous submission}
\institute{Affiliations withheld for review}

\maketitle

\begin{abstract}
Route guidance systems normally commit to one routing objective before any trip is
planned, even though the relative merit of established objectives can depend on the
traffic state. This paper studies the selection of the routing \emph{policy} rather than
of a route. Using 432 completed SUMO simulations over 24 controlled contexts on a
twelve-signal urban mesh, we evaluate shortest-distance routing (SP), dynamic
travel-time routing (DTT) and constrained travel-time-aware eco-routing under three
demand levels, four restriction configurations and two signal plans, with preference
expressed as a fixed-reference additive cost over mean journey time and tailpipe \CO.
SP has the lowest observed cost in all eight low-demand contexts and DTT in all eight
high-demand contexts, and the signal plan displaces the crossing by at least one step of
the demand grid. Relative to the single best fixed policy, leave-one-context-out
selection lowers mean decision regret from 0.03018 to 0.00606 normalised cost units ---
79.9\,\% of the available decision headroom captured, with the lowest-cost policy chosen
in 21 of 24 contexts. A two-threshold rule on demand and signal plan reproduces that
result exactly, so the benefit follows from the regime structure rather than from model
capacity. A leave-one-demand-out test reverses the conclusion: pooled regret rises to
0.07649, above the fixed-policy benchmark. Policy selection therefore has measurable
value inside the regimes represented during fitting and can be worse than a fixed policy
outside them. The contribution is the empirical characterisation of this decision
boundary and its limits; no new routing, multi-criteria or learning algorithm is
proposed.
\keywords{Routing-policy selection \and Urban traffic simulation \and Algorithm
selection \and Decision regret \and Traffic regimes \and Generalisation.}
\end{abstract}

\section{Introduction}

A route guidance system makes two decisions that are easy to conflate. The first is
the classical route-choice problem: which path should a vehicle take from an origin to
a destination? The second is rarely treated as a decision in its own right: which
routing \emph{policy} should generate that path? An operator may minimise static
distance, respond to recently observed travel times, or minimise estimated emissions
subject to a travel-time budget. Each policy returns a different path for the same
request, so the choice of policy precedes and constrains the choice of route. This
paper studies the second decision.

The distinction matters because congestion changes the relative value of routing
objectives. On a lightly loaded network the shortest path is usually also fast, and
minimising distance gives up little. As load rises, repeated concentration of traffic
on the shortest corridor raises delay there faster than on the alternatives, and a
policy that reacts to observed travel times can become preferable. Signal timing
modifies the same relationship, because the green split determines how much capacity
each competing movement receives. These mechanisms are individually well documented at
the level of routes and traffic dynamics, but they do not answer the question an
operator faces: should one routing objective be fixed for the whole day, or should the
active objective change with the operating context, and how much is that change worth?

We treat this as a context-dependent algorithm-selection problem in the sense of
Rice~\cite{rice1976algorithm}. The object being selected is a routing policy, not a
route: each policy remains responsible for assigning paths according to its own
objective, and a context-aware layer only decides which policy is active for the
current operating episode. Separating the two levels makes it possible to measure the
value of adapting the policy without attributing it to any particular route-generation
mechanism.

The central research question is:

\begin{quote}
\emph{Under which traffic regimes does the preferred established routing policy
change, how much decision value is available from adapting that policy, and how
reliably can the resulting structure be identified on unseen contexts?}
\end{quote}

Three subsidiary questions follow. How does the signal plan move the demand range in
which the preferred policy changes? How much of the available decision headroom can a
simple interpretable rule recover on held-out contexts, and does a flexible model
recover more? And what happens when the rule is applied outside the range of operating
conditions represented while it was fitted?

The paper makes four empirical contributions on a benchmark of 432 completed
simulations. It characterises a regime-dependent preference structure among established
routing policies on a signalised network, and shows that the signal plan displaces the
switching range by at least one step of the tested demand grid. It quantifies the value
of adaptation with context-level decision regret against a hindsight benchmark rather
than with prediction error, and reports how that value is distributed across regimes. It
shows that a two-threshold rule on demand and signal plan attains the same held-out
decision quality as a fitted regression tree, which bounds what the learned model
contributes. And it identifies an explicit generalisation boundary: the procedure that
helps inside the represented regime space becomes worse than the fixed policy it
replaces when an entire demand regime lies outside it. No new routing algorithm,
multi-criteria method or learning method is proposed, and no claim of priority is
made.
